% TOP_PROBLEM: {"id": 3102, "title": "Pebbling a cartesian product", "category_id": 3, "category": {"id": 3, "name": "graph_theory", "display_name": "Graph Theory", "description": "Problems involving graphs, networks, and their properties.", "slug": "graph-theory", "order_index": 3, "created_at": "2026-07-31T15:26:25.671Z"}, "status": "open", "problem_number": "OPG-586", "set_id": 12}
% FIRST_POSTED: 2026-10-01
% ENTRY_KIND: statement_only
% UnsolvedMath ID 3102; source catalogue code OPG-586
% !TeX program = lualatex
% Definition and sources only; this file is not an LLM solution attempt.
\documentclass[11pt,a4paper]{article}
\usepackage[margin=25mm]{geometry}
\usepackage{fontspec}
\setmainfont{lmroman10-regular.otf}[BoldFont=lmroman10-bold.otf,
 ItalicFont=lmroman10-italic.otf,BoldItalicFont=lmroman10-bolditalic.otf]
\usepackage{amsmath,amssymb,mathtools}
\usepackage{unicode-math}
\setmathfont{latinmodern-math.otf}
\AtBeginDocument{\let\setminus\smallsetminus}
\usepackage{xurl,hyperref}
\hypersetup{colorlinks=true,urlcolor=blue}
\setcounter{secnumdepth}{0}
\setlength{\parindent}{0pt}
\setlength{\parskip}{5pt}
\setlength{\emergencystretch}{3em}
\begin{document}
\section*{Pebbling a cartesian product}\label{3102}
{\small UnsolvedMath ID 3102 · OPG-586 · Graph Theory · OpenGarden}
\subsection{Definitions and mathematical statement}
Let $G$ be a finite nonempty connected simple undirected graph.
A pebble configuration is a function $C:V(G)\to\mathbb Z_{\ge0}$,
with total size $\sum_v C(v)$. A legal pebbling move chooses adjacent
vertices $u,v$ with $C(u)\ge2$, removes two pebbles from $u$, and adds
one pebble at $v$. One pebble is therefore lost at each move.

For a chosen target vertex $t$, a configuration is solvable at $t$ if
some finite sequence of legal moves leaves at least one pebble there.
The pebbling number $p(G)$ is the least integer $N$ such that every
configuration of total size $N$ is solvable at every target vertex.
The sequence of moves may depend on both the configuration and the
target; simultaneous delivery to all targets is not required.

For two such graphs, their Cartesian product $G_1\mathbin\Box G_2$ has
vertex set $V(G_1)\times V(G_2)$. Two distinct pairs are adjacent if
one coordinate is equal and the other coordinate is an adjacent pair
in its factor graph.
Graham's conjecture asserts
\[
 p(G_1\mathbin\Box G_2)\le p(G_1)p(G_2)
\]
for every pair of finite connected simple graphs.

Thus the product bound must handle arbitrary, possibly highly uneven,
placements on the product graph. A construction that succeeds only for
placements distributed uniformly among factor fibres is insufficient.
Connectedness ensures that pebbling numbers are finite; the single-vertex
graph is permitted and has pebbling number one.
\subsection{Short English statement}
Is the pebbling number of a Cartesian product at most the product of the pebbling numbers of its factors?
\subsection{Sources}
\begin{itemize}
\item[S1] ULAM.AI and the UnsolvedMath Contributors, supplied problems.json, record 3102 (OPG-586), OpenGarden collection. Catalogue identity and source transcription; CC BY 4.0.
\url{https://huggingface.co/datasets/ulamai/UnsolvedMath}.
\item[S2] Open Problem Garden, Pebbling a cartesian product. Original problem and discussion; the mathematical formulation below is expanded from the supplied source record.
\url{http://www.openproblemgarden.org/op/pebbling_a_cartesian_product}.
\end{itemize}
\end{document}
